Acemoglu, D., Johnson, S. \& Robinson, J.A. (2001). The Colonial Origins of Comparative Development: An Empirical Investigation. \emph{American Economic Review}, 91(5), 1369--1401.

Acemoglu, D. \& Robinson, J.A. (2012). \emph{Why Nations Fail: The
Origins of Power, Prosperity, and Poverty}. Crown Publishers.

Aghion, P. \& Howitt, P. (1992). A Model of Growth Through Creative
Destruction. \emph{Econometrica}, 60(2), 323--351.

Anderson, B.A. \& Silver, B.D. (1986). Infant Mortality in the Soviet
Union: Regional Differences and Measurement Issues. \emph{Population
and Development Review}, 12(4), 705--738.

Angrist, N., Djankov, S., Goldberg, P.K. \& Patrinos, H.A. (2021).
Measuring Human Capital Using Global Learning Data. \emph{Nature}, 592,
403--408.

Auvert, B., Taljaard, D., Lagarde, E., Sobngwi-Tambekou, J., Sitta, R.
\& Puren, A. (2005). Randomized, Controlled Intervention Trial of Male
Circumcision for Reduction of HIV Infection Risk: The ANRS 1265 Trial.
\emph{PLOS Medicine}, 2(11), e298.

Babiarz, K.S., Ma, P., Miller, G. \& Song, S. (2018). The Limits and
Consequences of Population Policy: Evidence from China's Wan Xi Shao
Campaign. NBER Working Paper No.~25130.

Baker, D.P., Collins, J.M. \& Leon, J. (2008). Risk Factor or Social
Vaccine? The Historical Progression of the Role of Education in HIV and
AIDS Infection in Sub-Saharan Africa. \emph{Prospects}, 38(4), 467--486.

Becker, G.S. (1964). \emph{Human Capital: A Theoretical and Empirical
Analysis, with Special Reference to Education}. National Bureau of
Economic Research; Columbia University Press.

Bora, J.K., Saikia, N., Kebede, E.B. \& Lutz, W. (2022). Revisiting
the causes of fertility decline in Bangladesh: the relative importance
of female education and family planning programs. \emph{Asian Population
Studies}, 19(1), 81--104.

Boyd, R. \& Richerson, P.J. (1985). \emph{Culture and the Evolutionary
Process}. University of Chicago Press.

Brass, P.R. (1986). The Political Uses of Crisis: The Bihar Famine of
1966--1967. \emph{Journal of Asian Studies}, 45(2), 245--267.

Cai, Y. (2010). China's Below-Replacement Fertility: Government Policy
or Socioeconomic Development? \emph{Population and Development Review},
36(3), 419--440.

Caldwell, J.C. (1979). Education as a Factor in Mortality Decline:
An Examination of Nigerian Data. \emph{Population Studies}, 33(3),
395--413.

Cleland, J.G. \& Van Ginneken, J.K. (1988). Maternal education and
child survival in developing countries: the search for pathways of
influence. \emph{Social Science \& Medicine}, 27(12), 1357--1368.

Davis, M. (2001). \emph{Late Victorian Holocausts: El Niño Famines and
the Making of the Third World}. Verso.

De~Neve, J.-W., Fink, G., Subramanian, S.V., Moyo, S. \& Bor, J.
(2015). Length of Secondary Schooling and Risk of HIV Infection in
Botswana: Evidence from a Natural Experiment. \emph{Lancet Global
Health}, 3(8), e470--e477.

de Waal, F.B.M. (2013). \emph{The Bonobo and the Atheist: In Search
of Humanism Among the Primates}. W.W. Norton.

De~Walque, D. (2007). How Does the Impact of an HIV/AIDS Information
Campaign Vary with Educational Attainment? Evidence from Rural Uganda.
\emph{Journal of Development Economics}, 84(2), 686--714.

Deaton, A. (2013). \emph{The Great Escape: Health, Wealth, and the
Origins of Inequality}. Princeton University Press.

Dehaene, S. (1997). \emph{The Number Sense: How the Mind Creates
Mathematics}. Oxford University Press.

Dehaene, S. (2009). \emph{Reading in the Brain: The New Science of
How We Read}. Viking.

Dehaene, S., Pegado, F., Braga, L.W., Ventura, P., Nunes Filho, G.,
Jobert, A., Dehaene-Lambertz, G., Kolinsky, R., Morais, J. \& Cohen, L.
(2010). How Learning to Read Changes the Cortical Networks for Vision and
Language. \emph{Science}, 330(6009), 1359--1364.

Diamond, J. (1997). \emph{Guns, Germs, and Steel: The Fates of Human
Societies}. W.W. Norton.

Drèze, J. \& Sen, A. (1989). \emph{Hunger and Public Action}. Clarendon
Press.

Drèze, J. \& Murthi, M. (2001). Fertility, Education, and Development:
Evidence from India. \emph{Population and Development Review}, 27(1),
33--63.

Duflo, E. (2012). Women Empowerment and Economic Development.
\emph{Journal of Economic Literature}, 50(4), 1051--1079.

Dyson, T. \& Maharatna, A. (1992). Bihar Famine, 1966--67, and
Maharashtra Drought, 1970--73: The Demographic Consequences.
\emph{Economic and Political Weekly}, 27(26), 1325--1332.

Easterlin, R.A. (1981). Why Isn't the Whole World Developed?
\emph{Journal of Economic History}, 41(1), 1--19.

Eveleth, P.B. \& Tanner, J.M. (1990). \emph{Worldwide Variation in
Human Growth} (2nd ed.). Cambridge University Press.

Fortson, J.G. (2008). The Gradient in Sub-Saharan Africa: Socioeconomic
Status and HIV/AIDS. \emph{Demography}, 45(2), 303--322.

Frank, M.C., Everett, D.L., Fedorenko, E. \& Gibson, E. (2008). Number
as a cognitive technology: Evidence from Pirahã language and cognition.
\emph{Cognition}, 108(3), 819--824.

Gakidou, E., Cowling, K., Lozano, R. \& Murray, C.J.L. (2010).
Increased Educational Attainment and Its Effect on Child Mortality in
175 Countries between 1970 and 2009: A Systematic Analysis.
\emph{The Lancet}, 376(9745), 959--974.

Galor, O. (2011). \emph{Unified Growth Theory}. Princeton University
Press.

Gao, M. (1999). \emph{Gao Village: Rural Life in Modern China}.
University of Hawai'i Press.

Gao, M. (2008). \emph{The Battle for China's Past: Mao and the Cultural
Revolution}. Pluto Press.

Glaeser, E.L., La Porta, R., López-de-Silanes, F. \& Shleifer, A.
(2004). Do Institutions Cause Growth? \emph{Journal of Economic
Growth}, 9(3), 271--303.

Goldin, C. \& Katz, L.F. (2008). \emph{The Race Between Education and
Technology}. Harvard University Press.

Goodall, J. (1986). \emph{The Chimpanzees of Gombe: Patterns of
Behavior}. Belknap Press of Harvard University Press.

Goody, J. (1976). \emph{Production and Reproduction: A Comparative
Study of the Domestic Domain}. Cambridge University Press.

Gordon, P. (2004). Numerical Cognition Without Words: Evidence from
Amazonia. \emph{Science}, 306(5695), 496--499.

Grantham-McGregor, S., Cheung, Y.B., Cueto, S., Glewwe, P., Richter,
L., Strupp, B.\ \& the International Child Development Steering Group.
(2007). Developmental Potential in the First 5 Years for Children in
Developing Countries. \emph{The Lancet}, 369(9555), 60--70.

Gwatkin, D.R. (1979). Political Will and Family Planning: The
Implications of India's Emergency Experience. \emph{Population and
Development Review}, 5(1), 29--59.

Hajnal, J. (1965). European Marriage Patterns in Perspective. In
D.V.~Glass \& D.E.C.~Eversley (Eds.), \emph{Population in History:
Essays in Historical Demography} (pp.~101--143). Edward Arnold.

Hamilton, W.D. (1964). The Genetical Evolution of Social Behaviour, I
\& II. \emph{Journal of Theoretical Biology}, 7(1), 1--52.

Hanushek, E.A. \& Woessmann, L. (2008). The Role of Cognitive Skills in
Economic Development. \emph{Journal of Economic Literature}, 46(3),
607--668.

Hanushek, E.A. \& Woessmann, L. (2012). Do Better Schools Lead to More
Growth? Cognitive Skills, Economic Outcomes, and Causation.
\emph{Journal of Economic Growth}, 17(4), 267--321.

Hare, B. \& Woods, V. (2020). \emph{Survival of the Friendliest:
Understanding Our Origins and Rediscovering Our Common Humanity}.
Random House.

Hawkes, K. (2003). Grandmothers and the evolution of human
longevity. \emph{American Journal of Human Biology}, 15(3), 380--400.

Heckman, J.J. (2006). Skill Formation and the Economics of Investing
in Disadvantaged Children. \emph{Science}, 312(5782), 1900--1902.

Heckman, J.J. \& Cunha, F. (2007). The Technology of Skill Formation.
\emph{American Economic Review}, 97(2), 31--47.

Heckman, J.J. \& Mosso, S. (2014). The Economics of Human Development
and Social Mobility. \emph{Annual Review of Economics}, 6, 689--733.

Herculano-Houzel, S. (2016). \emph{The Human Advantage: A New
Understanding of How Our Brain Became Remarkable}. MIT Press.

Heyes, C. (2018). \emph{Cognitive Gadgets: The Cultural Evolution of
Thinking}. Harvard University Press.

Hill, K. \& Hurtado, A.M. (1996). \emph{Aché Life History: The Ecology
and Demography of a Foraging People}. Aldine de Gruyter.

Hrdy, S.B. (2009). \emph{Mothers and Others: The Evolutionary Origins
of Mutual Understanding}. Harvard University Press.

Jensen, R. (2010). The (Perceived) Returns to Education and the Demand
for Schooling. \emph{Quarterly Journal of Economics}, 125(2), 515--548.

Kaplan, H., Hill, K., Lancaster, J. \& Hurtado, A.M. (2000). A Theory
of Human Life History Evolution: Diet, Intelligence, and Longevity.
\emph{Evolutionary Anthropology}, 9(4), 156--185.

Konner, M. (2010). \emph{The Evolution of Childhood: Relationships,
Emotion, Mind}. Harvard University Press.

Kramer, K.L. (2011). The Evolution of Human Parental Care and
Recruitment of Juvenile Help. \emph{Trends in Ecology \& Evolution},
26(10), 533--540.

Liu, Y., Rao, K. \& Hsiao, W.C. (2003). Medical Expenditure and Rural
Impoverishment in China. \emph{Journal of Health, Population and
Nutrition}, 21(3), 216--222.

Lutz, W. (2013). Demographic Metabolism: A Predictive Theory of
Socioeconomic Change. \emph{Population and Development Review},
38(s1), 283--301.

Lutz, W., Muttarak, R. \& Striessnig, E. (2014). Universal Education
is Key to Enhanced Climate Adaptation. \emph{Science}, 346(6213),
1061--1062.

Lutz, W. \& Kebede, E. (2018). Education and Health: Redrawing the
Preston Curve. \emph{Population and Development Review}, 44(2),
343--361.

Lutz, W., Reiter, C., Özdemir, C., Yildiz, D., Guimaraes, R. \& Goujon,
A. (2021). Skills-adjusted Human Capital Shows Rising Global Gap.
\emph{Proceedings of the National Academy of Sciences}, 118(7),
e2015826118.

Malthus, T.R. (1798). \emph{An Essay on the Principle of Population}.
London: J.~Johnson.

Mankiw, N.G., D. Romer, \& Weil, D.N. (1992). A Contribution to the
Empirics of Economic Growth. \emph{Quarterly Journal of Economics},
107(2), 407--437.

Marlowe, F.W. (2010). \emph{The Hadza: Hunter-Gatherers of Tanzania}.
University of California Press.

Maynard Smith, J. \& Szathmáry, E. (1995). \emph{The Major Transitions
in Evolution}. Oxford University Press.

Mokyr, J. (1983). \emph{Why Ireland Starved: A Quantitative and
Analytical History of the Irish Economy, 1800--1850}. George Allen
\& Unwin.

Mokyr, J. (1990). \emph{The Lever of Riches: Technological Creativity
and Economic Progress}. Oxford University Press.

Mokyr, J. (2002). \emph{The Gifts of Athena: Historical Origins of the
Knowledge Economy}. Princeton University Press.

Mokyr, J. (2009). \emph{The Enlightened Economy: An Economic History
of Britain, 1700--1850}. Yale University Press.

Mokyr, J. (2016). \emph{A Culture of Growth: The Origins of the Modern
Economy}. Princeton University Press.

Morris, I. (2010). \emph{Why the West Rules---For Now: The Patterns of History, and What They Reveal About the Future}. Farrar, Straus and Giroux.

Mullainathan, S. \& Shafir, E. (2013). \emph{Scarcity: Why Having Too
Little Means So Much}. Times Books.

Myhrvold-Hanssen, T. (2003). Democracy, News Media, and Famine
Prevention: Amartya Sen and the Bihar Famine of 1966--67. Masters thesis,
University of Oslo.

National Bureau of Statistics of China. (1982). \emph{Third National
Population Census of the People's Republic of China}. China Statistics
Press. {[}Rural population share: 79.7\% in 1980.{]}

Ó~Gráda, C. (1995). \emph{Ireland: A New Economic History
1780--1939}. Oxford University Press.

Ó~Gráda, C. (1999). \emph{Black '47 and Beyond: The Great Irish
Famine in History, Economy, and Memory}. Princeton University Press.

Pearl, J. (2009). \emph{Causality: Models, Reasoning, and Inference}
(2nd ed.). Cambridge University Press.

Pepper, S. (1996). \emph{Radicalism and Education Reform in 20th-Century
China}. Cambridge University Press.

Pinker, S. (2018). \emph{Enlightenment Now: The Case for Reason, Science, Humanism, and Progress}. Viking.

Prieto, A. (1981). Cuba's National Literacy Campaign. \emph{Journal of
Reading}, 25(3), 215--221. {[}268,000 brigadistas mobilised for the 1961
campaign.{]}

Pritchett, L. (2013). \emph{The Rebirth of Education: Schooling Ain't
Learning}. Center for Global Development.

Romer, P.M. (1986). Increasing Returns and Long-Run Growth.
\emph{Journal of Political Economy}, 94(5), 1002--1037.

Romer, P.M. (1990). Endogenous Technological Change. \emph{Journal of
Political Economy}, 98(5), S71--S102.

Roser, M., Ortiz-Ospina, E. \& Ritchie, H. (2013, ongoing). Life
Expectancy. \emph{Our World in Data}.
\url{https://ourworldindata.org/life-expectancy}.

Rosling, H., Rosling, O. \& Rosling R\"{o}nnlund, A. (2018).
\emph{Factfulness: Ten Reasons We're Wrong About the World --- and
Why Things Are Better Than You Think}. Flatiron Books.

Schultz, T.W. (1961). Investment in Human Capital.
\emph{American Economic Review}, 51(1), 1--17.

Schultz, T.P. (2002). Why Governments Should Invest More to Educate
Girls. \emph{World Development}, 30(2), 207--225.

Sen, A. (1999). \emph{Development as Freedom}. Oxford University Press.

Smil, V. (2017). \emph{Energy and Civilization: A History}. MIT Press.

Smith, A. (1776). \emph{An Inquiry into the Nature and Causes of the
Wealth of Nations}.

Solow, R.M. (1956). A Contribution to the Theory of Economic Growth.
\emph{Quarterly Journal of Economics}, 70(1), 65--94.

Stearns, S.C. (1992). \emph{The Evolution of Life Histories}. Oxford
University Press.

Tainter, J.A. (1988). \emph{The Collapse of Complex Societies}. Cambridge University Press.

Tomasello, M. (2014). \emph{A Natural History of Human Thinking}.
Harvard University Press.

Tomasello, M. (2016). \emph{A Natural History of Human Morality}.
Harvard University Press.

Tooby, J. \& Cosmides, L. (2010). Groups in Mind: The Coalitional Roots
of War and Morality. In H. H\"{o}gh-Olesen (Ed.), \emph{Human Morality
and Sociality: Evolutionary and Comparative Perspectives} (pp.~191--234).
Palgrave Macmillan.

UNDP. (1990). \emph{Human Development Report 1990}. Oxford University
Press.

UNESCO. (2024). \emph{Global Education Monitoring Report 2024/5:
Leadership in Education --- Lead for Learning}. Paris: UNESCO.

Unger, J. (1982). \emph{Education Under Mao: Class and Competition in
Canton Schools, 1960--1980}. Columbia University Press.

United Nations. (2015). \emph{Transforming Our World: The 2030 Agenda
for Sustainable Development} (A/RES/70/1). New York: United Nations.

Vicziany, M. (1982). Coercion in a Soft State: The Family-Planning
Program of India. Part~I: The Myth of Voluntarism. \emph{Pacific
Affairs}, 55(3), 373--402.

Walker, R., Gurven, M., Hill, K., Migliano, A., Chagnon, N., De
Souza, R., Djurovic, G., Hames, R., Hurtado, A.M., Kaplan, H.,
Kramer, K., Oliver, W.J., Valeggia, C. \& Yamauchi, T. (2006). Growth
Rates and Life Histories in Twenty-Two Small-Scale Societies.
\emph{American Journal of Human Biology}, 18(3), 295--311.

Wallace, A.R. (1870). \emph{Contributions to the Theory of Natural
Selection: A Series of Essays}. Macmillan.

World Bank. World Development Indicators. Various years.

Wrangham, R. \& Peterson, D. (1996). \emph{Demonic Males: Apes and
the Origins of Human Violence}. Houghton Mifflin.

Wrangham, R. (2009). \emph{Catching Fire: How Cooking Made Us Human}.
Basic Books.

Wrigley, E.A. (2010). \emph{Energy and the English Industrial
Revolution}. Cambridge University Press.
